\subsection{Cones; contraction of a subspace}

Let X and Y be nonempty topological spaces, and let \(f: X \to Y\) be a continuous map. Let \(\operatorname{Côn}(f)\) be the cone of the map \(f\) and let \(s\) be its vertex. Denote \(\alpha_{f}^{\prime}: X \times I \to \operatorname{Côn}(f)\) and \(\beta_{f}^{\prime}: Y \to \operatorname{Côn}(f)\) the canonical maps. The restriction of \(\alpha_{f}^{\prime}\) to the subspace \(X \times \{0\}\) of \(X \times I\) is the constant map with image \(\{s\}\). The map \(\beta_{f}^{\prime}\) induces a homeomorphism from Y onto the base of the cone \(\operatorname{Côn}(f)\), by which we will identify these two spaces. Let us also note

\[
\sigma_ {f} ^ {\prime} \colon (\operatorname{Côn} (f) - \{s \}) \times \mathbf {I} \to \operatorname{Côn} (f) - \{s \}
\]

the canonical contraction and \(\rho_{f}^{\prime}\colon\operatorname{Côn}(f) - \{s\}\to\mathrm{Y}\) the canonical retraction of the cone with its vertex removed onto its base.

Let us set \(J = \pi_{0}(X)\) ; for every element j of J, denote by \(X_{j}\) the j-th component of X, i.e., let \(b_{j}\) be a point of \(X_{j}\) and denote by \(\gamma_{j}\) the class of the path \(t \mapsto \alpha_{f}^{\prime}(b_{j}, t)\) on \(\operatorname{Côn}(f)\) which connects \(s\) to \(f(b_{j})\) .

Let I be the image of the map \(\pi_{0}(f)\) : this is the set of path-connected components of Y that meet \(f(\mathrm{X})\) ; let us denote by \(\varphi\colon J\to I\) the map deduced from f by passing to path-connected components. For every element \(i\in I\) , denote by \(Y_{i}\) the component i of Y and choose a point \(a_{i}\) on \(Y_{i}\) .

For every element \(j\) of J, choose a path \(\mathrm{B}_{j}\) connecting the point \(f(b_{j})\) to the point \(a_{\varphi(j)}\) on \(\mathrm{Y}_{\varphi(j)}\) and denote by \(\beta_{j}\) its class. Denote by \(\psi_{j}\) the homomorphism from \(\pi_{1}(\mathrm{X},b_{j})\) into \(\pi_{1}(\mathrm{Y},a_{\varphi(j)})\) defined by

\[
\psi_ {j} (v) = \beta_ {j} ^ {- 1} f _ {*} (v) \beta_ {j}
\]

for \(v \in \pi_1(X, a_j)\) .

Let \(\sigma\colon I\to J\) be a section of the map \(\varphi\) . Set \(T=\sigma(I)\) and \(\tau=\sigma\circ\varphi\) ; observe that \(\varphi\circ\tau=\varphi\) .

PROPOSITION 4. --- Suppose that the path-connected components of Y are open. For every \(i \in I\) , let \(G_i\) be the quotient of the group \(\pi_1(Y_i, b_i)\) by the smallest normal subgroup containing the image of the homomorphisms \(\psi_j\) , for \(j \in \overline{\varphi}^{-1}(i)\) ; let us denote by \(p_i\) the canonical surjection from \(\pi_1(Y_i, b_i)\) onto \(G_i\) .

There exists a unique group homomorphism

\[
\mathrm{P} \colon \underset {i \in \mathrm{I}} {*} \mathrm{G} _ {i} * \mathrm{F} (\mathrm{J} - \mathrm{T}) \to \pi_ {1} (\operatorname{C} \hat {\mathrm{n}} (f), s)
\]

such that

\[
\begin{array}{c c} \mathsf {P} (p _ {i} (v)) = \operatorname{Int} (\gamma_ {\sigma (i)} \beta_ {\sigma (i)}) (v) & \text {for } i\in I \text { and } v\in\pi_{1} (Y_{i},b_{i}), \\ \mathsf {P} (j) = \gamma_ {j} \beta_ {j} \beta_ {\tau (j)} ^ {- 1} \gamma_ {\tau (j)} ^ {- 1} & \text {for } j\in J - T. \end{array}
\]

The homomorphism P is an isomorphism.

Let Y' denote the union of the path-connected components of Y that meet \(f(\mathrm{X})\) and let \(f'\colon X \to Y'\) be the map given by \(x \mapsto f(x)\) . The set Y' is an open subset of Y; its path-connected components are open. The cone \(\operatorname{Côn}(f')\) is identified with the path-connected component of s in \(\operatorname{Côn}(f)\) . This allows us to assume that \(Y = Y'\) , in other words that the map \(\pi_{0}(f)\) is surjective and \(I = \pi_{0}(Y)\) .

For every \(j \in \mathrm{J}\) , set \(\mathrm{V}_j = \alpha_f'(X_j \times ]0,1[)\) . By passing to subspaces, the map \(\alpha_f'\) induces a homeomorphism from \(X \times ]0,1[\) onto the complement of \(Y \cup \{s\}\) in \(\operatorname{Côn}(f)\) . Consequently, the sets \(V_j\) are the path-connected components of \(\operatorname{Côn}(f) - (Y \cup \{s\})\) .

For every \(i \in \mathrm{I}\) , set \(\mathrm{U}_i = (\rho_f')^{-1}(\mathrm{Y}_i)\) ; it is an open subset of \(\operatorname{Côn}(f)\) since \(\mathrm{Y}_i\) is open in \(\mathrm{Y}\) by hypothesis. For every \(j \in \overline{\varphi}^{-1}(i)\) , we have \(f(\mathrm{X}_j) \subset \mathrm{Y}_i\) and

\[
\mathrm{V} _ {j} \cup \mathrm{Y} _ {i} = \alpha_ {f} ^ {\prime} (\mathrm{X} _ {j} \times ] 0, 1 ]) \cup \mathrm{Y} _ {i},
\]

so that \(V_{j} \cup Y_{i}\) is a path-connected subset of \(\operatorname{Côn}(f)\) containing \(Y_{i}\) . Since \(U_{i}\) is the union of \(Y_{i}\) and of the sets \(V_{j}\) , for \(j \in \overline{\varphi}^{-1}(i)\) , it follows that \(U_{i}\) is path-connected.

Finally, the set \(\mathrm{C}'(\mathrm{X}) = \mathrm{C}\hat{\mathrm{o}}\mathrm{n}(f) - \mathrm{Y}\) is an open subset of \(\mathrm{C}\hat{\mathrm{o}}\mathrm{n}(f)\) ; it is contractible to s, hence path-connected.

The set C'(X) and the sets \(U_{i}\) , for \(i \in I\) , constitute a covering of \(\operatorname{Côn}(f)\) by open path-connected subsets, a covering to which we shall apply prop. 1 of IV, p. 412. Let I' be the set obtained by adjoining s to I; one endows it with a total order for which s is its least element.

For distinct elements \(i, i'\) of I, one has \(\mathrm{U}_{i} \cap \mathrm{U}_{i'} = \varnothing\) . For \(i \in \mathrm{I}\) , \(\mathrm{C}'(\mathrm{X}) \cap \mathrm{U}_{i}\) is the union of the sets \(\mathrm{V}_{j}\) , for \(j \in \overline{\varphi}^{-1}(i)\) ; they are path-connected and pairwise disjoint. The intersection of any three distinct sets of this covering is empty.

The frame \(\Gamma\) of the covering under consideration has as vertices the set I'. Its arrows are the triples \((s, i, V_j)\) , for \(j \in \mathbf{J}\) and \(i = \varphi(j)\) ; thus we identify the set of arrows of \(\Gamma\) with the set J.

For \(i \in I\) we choose as base point \(\mathsf{a}(i) = a_i \in \mathrm{U}_i\) ; we also set \(\mathsf{a}(s) = s \in \mathrm{C}'(\mathrm{X})\) .

For \(j \in J\) we set \(\mathsf{b}(j) = \alpha_f'(b_j, \frac{1}{2})\) . We denote \(B_1(j)\) the path in \(C'(X)\) with source \(\mathsf{b}(j)\) and endpoint \(\mathsf{a}(s)\) given by \(t \mapsto \alpha_{f}^{\prime}(b_{j}, (1 - t)/2)\) . We denote \(\mathrm{B}_{2}(j)\) the path in \(\mathrm{U}_{\varphi(j)}\) with source \(\mathfrak{b}(j)\) and endpoint \(\mathfrak{a}(\varphi(j)) = a_{\varphi(j)}\) , given by the juxtaposition of the path \(t \mapsto \alpha_{f}^{\prime}(b_{j}, (1 + t)/2)\) and the path \(\mathrm{B}_{j}\) . Then, the class of the path \(\mathrm{B}(j) = \overline{\mathrm{B}}_{1}(j) * \mathrm{B}_{2}(j)\) is equal to \(\gamma_{j}\beta_{j}\) .

We choose \(i_{0}=s\) .

We take as maximal oriented tree T the unique oriented tree of \(\Gamma\) whose set of arrows is \(\sigma(\mathrm{I})\) . We have \(\delta(s) = e\) , while for \(i \in \mathrm{I}\) , \(\delta(i) = [\mathrm{B}(\sigma(i))] = \gamma_{\sigma(i)}\beta_{\sigma(i)}\) .

Let \(j \in \mathrm{J}\) and let \(i = \varphi(j)\) .

The homomorphism \(\varphi_{j}\) from \(\pi_1(V_j, \mathfrak{b}(j))\) into \(\pi_1(C'(X), s)\) is the trivial homomorphism because \(C'(X)\) is contractible to s .

The map \(\alpha_{f}^{\prime}\) induces a homeomorphism of \(\mathrm{X}_{j} \times ]0,1[\) onto \(\mathrm{V}_{j}\) ; this homeomorphism induces an isomorphism from the group \(\pi_{1}(\mathrm{V}_{j},\mathsf{b}(j))\) onto the group \(\pi_{1}(\mathrm{X}_{j},b_{j}) = \pi_{1}(\mathrm{X},b_{j})\) . By passing to subspaces, the map \(\sigma_{f}^{\prime}\) induces a strong contraction of \(\mathrm{U}_{i}\) onto \(\mathrm{Y}_{i}\) , which induces an isomorphism from the group \(\pi_{1}(\mathrm{U}_{i},\mathsf{a}(i))\) onto the group \(\pi_{1}(\mathrm{Y},a_{i})\) . By these isomorphisms, the homomorphism

\[
\psi_ {j} \colon \pi_ {1} (\mathrm{V} _ {j}, \mathsf {b} (j)) \to \pi_ {1} (\mathrm{U} _ {\varphi (j)}, \mathsf {a} (\varphi (j)))
\]

is identified with the homomorphism \(\operatorname{Int}(\beta_{\varphi(j)}^{-1}) \circ f_*\) from \(\pi_1(\mathrm{X}, b_j)\) into \(\pi_1(\mathrm{Y}, a_i)\) .

Since C'(X) is contractible to s, the homomorphism \(\mu_{s}\) is the trivial homomorphism.

Let \(i \in I\) . The homomorphism \(\mu_i\) from \(\pi_1(U_i, \mathfrak{a}(i))\) into \(\pi_1(\operatorname{Côn}(f), s)\) is identified with the homomorphism from \(\pi_1(Y, a_i)\) into \(\pi_1(\operatorname{Côn}(f), s)\) composed of the homomorphism \(\operatorname{Int}(\delta(i))\) and the homomorphism from \(\pi_1(U_i, s)\) into \(\pi_1(\operatorname{Côn}(f), s)\) deduced from the canonical injection of \(U_i\) into \(\operatorname{Côn}(f) \) .

Finally, the homomorphism \(\mu\colon\mathrm{F}(\mathrm{J})\to\pi_{1}(\mathrm{C}\hat{\mathrm{o}}\mathrm{n}(f),s)\) is given by

\[
\mu (j) = \gamma_ {j} \beta_ {j} \beta_ {\tau (j)} ^ {- 1} \gamma_ {\tau (j)} ^ {- 1}.
\]

Let P' be the unique group homomorphism

\[
\mathrm{P} ^ {\prime} \colon \underset {i \in \mathrm{I}} {*} \pi_ {1} (\mathrm{Y} _ {i}, a _ {i}) * \mathrm{F} (\mathrm{J}) \to \pi_ {1} (\mathrm{C} \hat {\mathrm{o}} \mathrm{n} (f), s)
\]

which coincides with the homomorphism \(\operatorname{Int}(\delta(i)^{-1})\) on \(\pi_1(Y_i, a_i)\) and with the homomorphism \(\mu\) on \(\mathrm{F}(\mathrm{J})\). By prop. 1 of IV, p. 412, the homomorphism P' is surjective and its kernel is the smallest distinguished subgroup of \(*\pi_1(Y_i, a_i) * F(J)\) which contains the elements \(r_1(j)\) , for \(j \in T\) , and the elements \(r_{2}(j, v)\) , for \(j \in J\) and \(v \in \pi_{1}(V_{j}, b(j))\) . (There are no elements \(r_{3}(k)\) , because the set K is empty.)

For \(j \in \mathrm{T}\) , we have \(r_1(j) = j\) . Let \(j \in \mathrm{J}\) and let \(v \in \pi_1(\mathrm{V}_j, \mathfrak{b}(j))\) . Taking into account the identification of \(\pi_1(\mathrm{V}_j, \mathfrak{b}(j))\) with \(\pi_1(\mathrm{X}, b_j)\) , we have \(r_2(j, v) = j\psi_j(v)j^{-1}\) . Denote by \(p\) the canonical surjective homomorphism from \(*_{i}\pi_1(\mathrm{Y}_i, a_i) * \mathrm{F}(\mathrm{J})\) onto \(*_{i}\mathrm{G}_{i} * \mathrm{F}(\mathrm{J} - \mathrm{T})\) . For \(j \in \mathrm{T}\) , we have \(p(r_1(j)) = e\) ; for \(j \in \mathrm{J}\) and \(v \in \pi_1(\mathrm{X}, b_j)\) , we have \(p(\psi_j(v)) = e\) . Consequently, there exists a unique group homomorphism \(\mathsf{P}\) from \(*_{i}\mathrm{G}_{i} * \mathrm{F}(\mathrm{J} - \mathrm{T})\) into \(\pi_1(\operatorname{Côn}(f), s)\) such that \(\mathsf{P}' \circ p = \mathsf{P}\) ; it is an isomorphism.

COROLLARY 1. --- Suppose moreover that the spaces X and Y are path-connected and let a be a point of X. The canonical map from \(\pi_1(Y, f(a))\) into \(\pi_1(\operatorname{Côn}(f), f(a))\) is surjective, and its kernel is the smallest normal subgroup containing the image of the homomorphism \(\pi_1(f, a)\) .

COROLLARY 2. --- Suppose moreover that the path-connected components of Y are simply connected by arcs. The homomorphism \(\mu\colon F(J-T)\to\pi_{1}(\operatorname{Côn}(f),s)\) is then an isomorphism.

Remark 1. --- Let X be a topological space whose path-connected components are open. Let A be a closed subspace of X; denote by \(\iota\colon A\to X\) the canonical injection, by X/A the space obtained from X by contracting A to a point, and denote by \(o\) the point to which A is contracted and by \(p\colon X\to X / A\) the canonical map. Suppose in addition that the pair (X,A) has the homotopy extension property. The canonical map \(\overline{\rho}\colon \operatorname{Côn}(\iota)\to \mathrm{X} / \mathrm{A}\) is then a homotopy equivalence (III, p. 255, remark 1), and from prop. 4 one deduces the computation of the Poincaré group of X/A at its base point \(o\) . In particular, if the path-connected components of X are simply connected by arcs, the group \(\pi_1(\mathrm{X / A},o)\) is a free group.
% semantic-fragments: legacy-input-seam
\relax{}
